\section{Реализация}

\subsection{Структура приложения}

Приложение состоит из следующих классов и структур:

\begin{itemize}
    \item \texttt{Contact}~--- структура для хранения данных контакта;
    \item \texttt{PhoneNumber}~--- структура для хранения телефонного номера с типом;
    \item \texttt{MainWindow}~--- главное окно приложения, наследник \texttt{QMainWindow};
    \item \texttt{ContactDialog}~--- диалог создания и редактирования контакта, наследник \texttt{QDialog};
    \item \texttt{Validator}~--- класс с статическими методами валидации данных;
    \item \texttt{PhoneBookFile}~--- класс для работы с файлами, наследник \texttt{QFile}.
\end{itemize}

\subsection{Структуры данных}

\subsubsection*{Структура PhoneNumber}

Структура хранит телефонный номер и его тип. Тип определяется перечислением с вариантами: рабочий, домашний, мобильный.

\begin{lstlisting}[language=C, caption=Структура PhoneNumber]
структура PhoneNumber:
    type    // тип номера (Work, Home, Mobile)
    number  // номер как последовательность цифр
    
    функция typeString():
        вернуть строковое представление типа
    конец функции
    
    функция formatted():
        если длина(number) = 11 и number начинается с "7"
            вернуть "+X (XXX) XXX-XX-XX"
        иначе
            вернуть "+" + number
    конец функции
конец структуры
\end{lstlisting}

\subsubsection*{Структура Contact}

Структура объединяет все данные контакта и предоставляет вспомогательные методы.

\begin{lstlisting}[language=C, caption=Структура Contact]
структура Contact:
    firstName    // имя
    lastName     // фамилия
    middleName   // отчество
    address      // адрес
    birthDate    // дата рождения
    email        // электронная почта
    phones       // список телефонов
    
    функция fullName():
        вернуть lastName + " " + firstName + " " + middleName
    конец функции
    
    функция primaryPhone():
        если phones пуст
            вернуть ""
        вернуть phones[0].formatted()
    конец функции
конец структуры
\end{lstlisting}

Для сериализации структур определены операторы потокового ввода-вывода, позволяющие использовать \texttt{QDataStream} для бинарного сохранения.

\subsubsection*{Структура Result}

Структура используется для возврата результата валидации из методов класса \texttt{Validator}. Содержит флаг успешности проверки, сообщение об ошибке и очищенное значение.

\begin{lstlisting}[language=C, caption=Структура Result]
структура Result:
    valid    // флаг корректности (true/false)
    error    // сообщение об ошибке (пустое, если valid = true)
    cleaned  // очищенное значение после обработки
    
    конструктор Result(v, e, c):
        valid := v
        error := e
        cleaned := c
    конец конструктора
конец структуры
\end{lstlisting}

Такой подход позволяет методам валидации одновременно сообщать об ошибке и возвращать нормализованное значение (например, строку без лишних пробелов или номер телефона в виде последовательности цифр).

\subsection{Класс Validator}

Класс содержит статические методы для проверки корректности вводимых данных. Каждый метод возвращает структуру \texttt{Result} с полями: флаг валидности, сообщение об ошибке, очищенное значение.

\subsubsection*{Интерфейс класса}

\begin{lstlisting}[language=C, caption=Заголовочный файл validator.h]
класс Validator:
    // Публичные статические методы валидации
    статический validateName(name: QString) -> Result
    статический validatePhone(phone: QString) -> Result
    статический validateEmail(email: QString) -> Result
    статический validateBirthDate(date: QDate) -> Result
    статический validateAddress(address: QString) -> Result

    // Приватные вспомогательные методы
    приватный статический trimSpaces(str: QString) -> QString
    приватный статический isLeapYear(year: int) -> bool
    приватный статический daysInMonth(month: int, year: int) -> int
конец класса
\end{lstlisting}

\begin{itemize}
    \item \texttt{validateName}~--- проверяет ФИО на соответствие требованиям (заглавная буква, допустимые символы, отсутствие дефиса в начале/конце);
    \item \texttt{validatePhone}~--- проверяет телефон и извлекает из него последовательность цифр;
    \item \texttt{validateEmail}~--- проверяет формат email и удаляет пробелы;
    \item \texttt{validateBirthDate}~--- проверяет корректность даты и что она в прошлом;
    \item \texttt{validateAddress}~--- проверяет, что адрес не пустой, и удаляет лишние пробелы;
    \item \texttt{trimSpaces}~--- удаляет начальные, конечные и дублирующиеся пробелы;
    \item \texttt{isLeapYear}~--- определяет, является ли год високосным;
    \item \texttt{daysInMonth}~--- возвращает количество дней в указанном месяце.
\end{itemize}

\subsubsection*{Валидация имени}

\begin{lstlisting}[language=C, caption=Функция validateName]
функция validateName(name):
    cleaned := убрать_пробелы(name)
    
    если cleaned пуст
        вернуть Result(false, "Поле не может быть пустым")
    
    если cleaned не начинается с заглавной буквы
        вернуть Result(false, "Должно начинаться с заглавной буквы")
    
    если cleaned начинается с "-" или заканчивается на "-"
        вернуть Result(false, "Не может начинаться/заканчиваться на дефис")
    
    если cleaned не соответствует шаблону допустимых символов
        вернуть Result(false, "Недопустимые символы")
    
    вернуть Result(true, "", cleaned)
конец функции
\end{lstlisting}

Для проверки заглавной буквы используется паттерн \verb|^\p{Lu}|, где \verb|\p{Lu}| соответствует любой заглавной букве Unicode. Для проверки допустимых символов~--- паттерн \verb|^[\p{L}\p{N}\- ]+$|, разрешающий буквы, цифры, дефис и пробел.

\subsubsection*{Валидация телефона}

\begin{lstlisting}[language=C, caption=Функция validatePhone]
функция validatePhone(phone):
    cleaned := убрать_пробелы(phone)
    
    если cleaned пуст
        вернуть Result(false, "Телефон не может быть пустым")
    
    если cleaned не соответствует шаблону формата телефона
        вернуть Result(false, "Недопустимые символы")
    
    // извлекаем только цифры
    digitsOnly := извлечь_цифры(cleaned)
    
    если длина(digitsOnly) < 7
        вернуть Result(false, "Минимум 7 цифр")
    
    если длина(digitsOnly) > 15
        вернуть Result(false, "Слишком длинный номер")
    
    вернуть Result(true, "", digitsOnly)
конец функции
\end{lstlisting}

Паттерн \verb|^[\+]?[\d\s\(\)\-]+$| допускает опциональный плюс в начале, цифры, пробелы, скобки и дефисы.

\subsubsection*{Валидация email}

\begin{lstlisting}[language=C, caption=Функция validateEmail]
функция validateEmail(email):
    // удаляем все пробелы
    cleaned := удалить_все_пробелы(email)
    
    если cleaned пуст
        вернуть Result(false, "Email не может быть пустым")
    
    если cleaned не соответствует шаблону email
        вернуть Result(false, "Неверный формат email")
    
    вернуть Result(true, "", cleaned)
конец функции
\end{lstlisting}

Паттерн email проверяет наличие имени пользователя, символа @, домена и доменной зоны минимум из двух символов.

\subsubsection*{Валидация даты рождения}

\begin{lstlisting}[language=C, caption=Функция validateBirthDate]
функция validateBirthDate(date):
    если date некорректна
        вернуть Result(false, "Некорректная дата")
    
    если date >= текущая_дата()
        вернуть Result(false, "Дата должна быть в прошлом")
    
    month := date.month()
    day := date.day()
    year := date.year()
    
    если month < 1 или month > 12
        вернуть Result(false, "Месяц от 1 до 12")
    
    maxDays := daysInMonth(month, year)
    если day < 1 или day > maxDays
        вернуть Result(false, "Неверное число дней")
    
    вернуть Result(true, "", date.toString())
конец функции

функция daysInMonth(month, year):
    если month в [1, 3, 5, 7, 8, 10, 12]
        вернуть 31
    если month в [4, 6, 9, 11]
        вернуть 30
    если month = 2
        если isLeapYear(year)
            вернуть 29
        вернуть 28
конец функции
\end{lstlisting}

\subsection{Класс MainWindow}

Главное окно наследуется от \texttt{QMainWindow} и использует макрос \texttt{Q\_OBJECT} для поддержки сигналов и слотов.

\subsubsection*{Интерфейс класса}

\begin{lstlisting}[language=C, caption=Заголовочный файл mainwindow.h]
класс MainWindow : наследует QMainWindow:
    Q_OBJECT
    
    // Конструктор и деструктор
    публичный MainWindow(parent: QWidget*)
    публичный ~MainWindow()

    // Слоты (обработчики событий)
    приватный слот addContact()
    приватный слот editContact()
    приватный слот removeContact()
    приватный слот searchContacts()
    приватный слот sortByColumn(column: int)
    приватный слот onSelectionChanged()
    приватный слот exportToYaml()
    приватный слот importFromYaml()

    // Приватные методы
    приватный setupUI()
    приватный setupMenuBar()
    приватный loadData()
    приватный saveData()
    приватный updateTable()
    приватный updateTable(contacts: QList<Contact>)

    // Поля данных
    приватный m_table: QTableWidget*
    приватный m_searchEdit: QLineEdit*
    приватный m_searchFieldCombo: QComboBox*
    приватный m_contacts: QList<Contact>
    приватный m_dataFile: QString
    приватный m_sortColumn: int
    приватный m_sortOrder: Qt::SortOrder
    приватный m_sortClickCount: int
конец класса
\end{lstlisting}

\begin{itemize}
    \item \texttt{MainWindow}~--- конструктор, инициализирует интерфейс и загружает данные;
    \item \texttt{\textasciitilde MainWindow}~--- деструктор, сохраняет данные при закрытии;
    \item \texttt{addContact}~--- открывает диалог создания нового контакта;
    \item \texttt{editContact}~--- открывает диалог редактирования выбранного контакта;
    \item \texttt{removeContact}~--- удаляет выбранный контакт с подтверждением;
    \item \texttt{searchContacts}~--- фильтрует таблицу по введённому запросу;
    \item \texttt{sortByColumn}~--- сортирует таблицу по указанному столбцу;
    \item \texttt{onSelectionChanged}~--- включает/отключает кнопки при выборе строки;
    \item \texttt{exportToYaml}~--- экспортирует контакты в YAML-файл;
    \item \texttt{importFromYaml}~--- импортирует контакты из YAML-файла;
    \item \texttt{setupUI}~--- создаёт виджеты и компоновку интерфейса;
    \item \texttt{setupMenuBar}~--- создаёт меню приложения;
    \item \texttt{loadData}~--- загружает контакты из файла при запуске;
    \item \texttt{saveData}~--- сохраняет контакты в файл;
    \item \texttt{updateTable}~--- обновляет содержимое таблицы.
\end{itemize}

\subsubsection*{Используемые виджеты Qt}

\begin{itemize}
    \item \texttt{QTableWidget}~--- таблица для отображения контактов;
    \item \texttt{QLineEdit}~--- поле ввода для поиска;
    \item \texttt{QComboBox}~--- выпадающий список для выбора поля поиска;
    \item \texttt{QPushButton}~--- кнопки управления;
    \item \texttt{QMenuBar}, \texttt{QMenu}, \texttt{QAction}~--- меню приложения.
\end{itemize}

\subsubsection*{Сигналы и слоты}

Механизм сигналов и слотов Qt используется для связывания действий пользователя с методами обработки:

\begin{lstlisting}[language=C, caption=Подключение сигналов и слотов]
// кнопка "Добавить" -> метод addContact
connect(m_addButton, clicked, this, addContact)

// кнопка "Редактировать" -> метод editContact
connect(m_editButton, clicked, this, editContact)

// кнопка "Удалить" -> метод removeContact
connect(m_removeButton, clicked, this, removeContact)

// изменение текста поиска -> метод searchContacts
connect(m_searchEdit, textChanged, this, searchContacts)

// клик по заголовку столбца -> метод sortByColumn
connect(m_table.header, sectionClicked, this, sortByColumn)

// двойной клик по строке -> редактирование
connect(m_table, doubleClicked, this, editContact)
\end{lstlisting}

\subsubsection*{Алгоритм поиска}

\begin{lstlisting}[language=C, caption=Функция searchContacts]
функция searchContacts():
    query := в_нижний_регистр(убрать_пробелы(m_searchEdit.text()))
    field := m_searchFieldCombo.currentData()
    
    если query пуст
        updateTable(m_contacts)
        выйти из функции
    
    filtered := пустой список
    
    для каждого contact в m_contacts
        match := false
        
        выбор field
            случай LAST_NAME:
                match := содержит(contact.lastName, query)
            случай FIRST_NAME:
                match := содержит(contact.firstName, query)
            // ... остальные поля аналогично
        конец выбора
        
        если match = true
            добавить contact в filtered
    конец цикла
    
    updateTable(filtered)
конец функции
\end{lstlisting}

\subsubsection*{Алгоритм сортировки}

Реализована трёхступенчатая сортировка: первый клик~--- по возрастанию, второй~--- по убыванию, третий~--- сброс к исходному порядку.

\begin{lstlisting}[language=C, caption=Функция sortByColumn]
функция sortByColumn(column):
    если column = m_sortColumn
        m_sortClickCount := m_sortClickCount + 1
        
        если m_sortClickCount = 1
            // второй клик - сортировка по убыванию
            m_sortOrder := DescendingOrder
            m_table.sortItems(column, m_sortOrder)
            показать_индикатор(column, m_sortOrder)
        иначе
            // третий клик - сброс сортировки
            m_sortColumn := -1
            m_sortClickCount := 0
            скрыть_индикатор()
            searchContacts()  // восстановить исходный порядок
    иначе
        // новый столбец - сортировка по возрастанию
        m_sortColumn := column
        m_sortClickCount := 0
        m_sortOrder := AscendingOrder
        m_table.sortItems(column, m_sortOrder)
        показать_индикатор(column, m_sortOrder)
конец функции
\end{lstlisting}

\subsection{Класс ContactDialog}

Диалог создания и редактирования контакта наследуется от \texttt{QDialog}.

\subsubsection*{Интерфейс класса}

\begin{lstlisting}[language=C, caption=Заголовочный файл createcontact.h]
класс ContactDialog : наследует QDialog:
    Q_OBJECT
    
    // Конструктор
    публичный ContactDialog(parent: QWidget*)

    // Публичные методы
    публичный setContact(contact: Contact)
    публичный getContact() -> Contact

    // Слоты
    приватный слот addPhone()
    приватный слот removePhone()
    приватный слот validate()

    // Приватные методы
    приватный setupUI()
    приватный updatePhoneTable()
    приватный validateAll() -> bool
    приватный showError(label: QLabel*, error: QString)
    приватный clearError(label: QLabel*)

    // Поля ввода
    приватный m_lastNameEdit: QLineEdit*
    приватный m_firstNameEdit: QLineEdit*
    приватный m_middleNameEdit: QLineEdit*
    приватный m_addressEdit: QLineEdit*
    приватный m_birthDateEdit: QDateEdit*
    приватный m_emailEdit: QLineEdit*
    приватный m_phonesTable: QTableWidget*
    приватный m_phoneInput: QLineEdit*
    приватный m_phoneTypeCombo: QComboBox*
    
    // Метки ошибок
    приватный m_lastNameError: QLabel*
    // ... остальные метки аналогично
    
    // Данные
    приватный m_phones: QList[PhoneNumber]
конец класса
\end{lstlisting}

\begin{itemize}
    \item \texttt{ContactDialog}~--- конструктор, создаёт интерфейс диалога;
    \item \texttt{setContact}~--- заполняет поля данными существующего контакта для редактирования;
    \item \texttt{getContact}~--- собирает данные из полей и возвращает объект контакта;
    \item \texttt{addPhone}~--- добавляет введённый телефон в список после валидации;
    \item \texttt{removePhone}~--- удаляет выбранный телефон из списка;
    \item \texttt{validate}~--- вызывает \texttt{validateAll} и закрывает диалог при успехе;
    \item \texttt{setupUI}~--- создаёт все виджеты и настраивает компоновку;
    \item \texttt{updatePhoneTable}~--- обновляет таблицу телефонов;
    \item \texttt{validateAll}~--- проверяет все поля формы и показывает ошибки;
    \item \texttt{showError}~--- отображает текст ошибки в метке;
    \item \texttt{clearError}~--- очищает метку ошибки.
\end{itemize}

\subsubsection*{Используемые виджеты}

\begin{itemize}
    \item \texttt{QLineEdit}~--- поля ввода для ФИО, адреса, email, телефона;
    \item \texttt{QDateEdit}~--- выбор даты с календарём;
    \item \texttt{QComboBox}~--- выбор типа телефона;
    \item \texttt{QTableWidget}~--- список добавленных телефонов;
    \item \texttt{QLabel}~--- отображение ошибок валидации.
\end{itemize}

\subsubsection*{Алгоритм валидации формы}

\begin{lstlisting}[language=C, caption=Функция validateAll]
функция validateAll():
    valid := true
    
    // проверка фамилии
    result := Validator.validateName(m_lastNameEdit.text())
    если result.valid = false
        showError(m_lastNameError, result.error)
        valid := false
    иначе
        clearError(m_lastNameError)
    
    // проверка имени
    result := Validator.validateName(m_firstNameEdit.text())
    если result.valid = false
        showError(m_firstNameError, result.error)
        valid := false
    иначе
        clearError(m_firstNameError)
    
    // отчество - необязательное поле
    если m_middleNameEdit.text() не пуст
        result := Validator.validateName(m_middleNameEdit.text())
        если result.valid = false
            showError(m_middleNameError, result.error)
            valid := false
    
    // проверка адреса, даты, email аналогично...
    
    // минимум один телефон
    если m_phones пуст
        showError(m_phoneError, "Добавьте хотя бы один телефон")
        valid := false
    
    вернуть valid
конец функции
\end{lstlisting}

\subsection{Класс PhoneBookFile}

Класс наследуется от \texttt{QFile} и расширяет его методами для работы с контактами.

\subsubsection*{Интерфейс класса}

\begin{lstlisting}[language=C, caption=Заголовочный файл phonebooktofile.h]
класс PhoneBookFile : наследует QFile:
    Q_OBJECT
    
    // Конструктор
    публичный PhoneBookFile(fileName: QString, parent: QObject*)

    // Методы работы с бинарным форматом
    публичный loadContacts(contacts: QList[Contact]) -> bool
    публичный saveContacts(contacts: QList[Contact]) -> bool
    
    // Методы работы с YAML
    публичный loadContactsYaml(contacts: QList[Contact]) -> bool
    публичный saveContactsYaml(contacts: QList[Contact]) -> bool
    
    // Получение ошибки
    публичный lastError() -> QString

    // Константы формата файла
    статический FILE_MAGIC: quint32 = 0x50484F4E  // "PHON"
    статический FILE_VERSION: quint32 = 1

    // Приватные поля
    приватный m_lastError: QString
конец класса
\end{lstlisting}

\begin{itemize}
    \item \texttt{PhoneBookFile}~--- конструктор, принимает имя файла;
    \item \texttt{loadContacts}~--- загружает контакты из бинарного файла;
    \item \texttt{saveContacts}~--- сохраняет контакты в бинарный файл;
    \item \texttt{loadContactsYaml}~--- загружает контакты из YAML-файла;
    \item \texttt{saveContactsYaml}~--- сохраняет контакты в YAML-файл;
    \item \texttt{lastError}~--- возвращает текст последней ошибки;
    \item \texttt{FILE\_MAGIC}~--- магическое число для идентификации формата файла;
    \item \texttt{FILE\_VERSION}~--- версия формата для обеспечения совместимости.
\end{itemize}

\subsubsection*{Бинарный формат}

Для бинарного хранения используется \texttt{QDataStream}. Файл начинается с <<магического>> числа и версии для проверки совместимости.

\begin{lstlisting}[language=C, caption=Функция saveContacts]
функция saveContacts(contacts):
    если не удалось открыть файл для записи
        m_lastError := "Не удалось открыть файл"
        вернуть false
    
    stream := новый QDataStream(this)
    
    // заголовок файла
    записать в stream значение FILE_MAGIC   // 0x50484F4E ("PHON")
    записать в stream значение FILE_VERSION // 1
    
    // данные
    записать в stream значение contacts
    
    закрыть файл
    вернуть true
конец функции
\end{lstlisting}

\begin{lstlisting}[language=C, caption=Функция loadContacts]
функция loadContacts(contacts):
    contacts.clear()
    
    если файл не существует
        вернуть true  // пустой список при первом запуске
    
    если не удалось открыть файл для чтения
        m_lastError := "Не удалось открыть файл"
        вернуть false
    
    stream := новый QDataStream(this)
    
    прочитать из stream в magic
    прочитать из stream в version
    
    если magic != FILE_MAGIC
        m_lastError := "Неверный формат файла"
        вернуть false
    
    если version > FILE_VERSION
        m_lastError := "Файл от новой версии программы"
        вернуть false
    
    прочитать из stream в contacts
    
    закрыть файл
    вернуть true
конец функции
\end{lstlisting}

\subsubsection*{Формат YAML}

Для работы с YAML используется библиотека yaml-cpp. Каждый контакт сохраняется как словарь с вложенным списком телефонов.

\begin{lstlisting}[language=C, caption=Функция saveContactsYaml]
функция saveContactsYaml(contacts):
    out := новый YAML::Emitter()
    
    начать последовательность в out
    
    для каждого contact в contacts
        начать словарь в out
        
        записать ключ "lastName" со значением contact.lastName
        записать ключ "firstName" со значением contact.firstName
        записать ключ "middleName" со значением contact.middleName
        записать ключ "address" со значением contact.address
        записать ключ "email" со значением contact.email
        записать ключ "birthDate" со значением contact.birthDate
        
        записать ключ "phones" и начать вложенную последовательность
        для каждого phone в contact.phones
            начать словарь
            записать ключ "type" со значением phone.typeString()
            записать ключ "number" со значением phone.number
            закончить словарь
        конец цикла
        закончить последовательность
        
        закончить словарь
    конец цикла
    
    закончить последовательность
    
    записать out.c_str() в файл
    вернуть true
конец функции
\end{lstlisting}

\begin{lstlisting}[language=C, caption=Функция loadContactsYaml]
функция loadContactsYaml(contacts):
    contacts.clear()
    
    root := YAML::LoadFile(fileName())
    
    если root не является списком
        m_lastError := "Ожидался список контактов"
        вернуть false
    
    для каждого node в root
        contact := новый Contact()
        
        если node["lastName"] существует
            contact.lastName := прочитать node["lastName"] как строку
        если node["firstName"] существует
            contact.firstName := прочитать node["firstName"] как строку
        // ... остальные поля аналогично
        
        если node["phones"] существует и является списком
            для каждого phoneNode в node["phones"]
                phone := новый PhoneNumber()
                phone.type := разобрать_тип(прочитать phoneNode["type"])
                phone.number := прочитать phoneNode["number"] как строку
                добавить phone в contact.phones
            конец цикла
        
        добавить contact в contacts
    конец цикла
    
    вернуть true
конец функции
\end{lstlisting}

\subsection{Наследование от классов Qt}

\subsubsection*{MainWindow от QMainWindow}

Наследование от \texttt{QMainWindow} обеспечивает:
\begin{itemize}
    \item готовую структуру главного окна с центральным виджетом;
    \item поддержку меню через \texttt{menuBar()};
    \item возможность добавления панелей инструментов и строки состояния.
\end{itemize}

\subsubsection*{ContactDialog от QDialog}

Наследование от \texttt{QDialog} предоставляет:
\begin{itemize}
    \item модальное поведение окна;
    \item методы \texttt{accept()} и \texttt{reject()} для закрытия с результатом;
    \item интеграцию с \texttt{exec()} для блокирующего вызова.
\end{itemize}

\subsubsection*{PhoneBookFile от QFile}

Наследование от \texttt{QFile} даёт:
\begin{itemize}
    \item готовые методы работы с файлами: \texttt{open()}, \texttt{close()}, \texttt{exists()};
    \item интеграцию с \texttt{QDataStream} для бинарной сериализации;
    \item обработку ошибок через \texttt{errorString()}.
\end{itemize}